% Coverage: physical pp. 66--71 (printed pp. 43--48), from the Section 25.4
% heading through the final paragraph on physical p. 70, ending immediately
% before Section 25.5 on physical p. 71. Equations (25.4.1)--(25.4.9), one
% unnumbered display, citation [3], three footnotes, and no figures, tables, or
% typographic dividers. Uncertainty: none.

\subsection{Massless Particle Supermultiplets}

\providecommand{\WeylQ}{\mathrm{Q}}

Supersymmetry requires that the known particles be accompanied in irreducible
representations of the supersymmetry algebra by `sparticles': bosonic
`squarks' and `sleptons' accompanying quarks and leptons, and fermionic
`gauginos' accompanying gauge bosons. None of these sparticles have been
observed, so supersymmetry is certainly broken, with the masses of the
sparticles almost certainly much larger than the quark, lepton, and gauge
boson masses produced by the spontaneous breakdown of the electroweak
$SU(2)\times U(1)$ gauge group, and hence of the same order of magnitude as the
splittings within supermultiplets. Thus it is very likely that at energy
scales that are large enough so that we can neglect supersymmetry breaking and
these mass splittings, we can also treat the known quarks, leptons, and gauge
bosons and their superpartners as massless. Therefore we will be specially
interested in supermultiplets of massless particles.

Consider a state containing a single massless particle belonging to some
supermultiplet. We obtain the other states in the same supermultiplet by
applying operators $\WeylQ_{ar}$ and/or $\WeylQ^\dagger_{ar}$ to this state.
Since $\WeylQ_{ar}$ and $\WeylQ^\dagger_{ar}$ commute with $P_\mu$, all these
states have the same value of the four-momentum. We will work in a Lorentz
frame in which the four-momentum of these states is $p^1=p^2=0$ and
$p^3=p^0=E$. With this choice of four-momentum, we have
\begin{equation}
  \sigma_\mu p^\mu
  =E(\sigma_0+\sigma_3)
  =2E
  \begin{pmatrix}
    1&0\\
    0&0
  \end{pmatrix},
  \tag{25.4.1}\label{eq:25.4.1}
\end{equation}
which aside from the factor $2E$ is the projection matrix onto the subspace
with helicity $+1/2$. The anticommutation relation \eqref{eq:25.2.7} therefore
shows that
$\{\WeylQ_{(-1/2)r},\WeylQ^\dagger_{(-1/2)r}\}$ gives zero when acting on
any state in a supermultiplet with this momentum, and thus so do
$\WeylQ_{(-1/2)r}$ and $\WeylQ^\dagger_{(-1/2)r}$. We must therefore
construct the states of the supermultiplet by acting only with
$\WeylQ_{(1/2)r}$ and $\WeylQ^\dagger_{(1/2)r}$. Furthermore, we are
labelling the $\WeylQ$s with their $J_3$ values, in the sense that
\begin{equation}
  [J_3,\WeylQ_{ar}]=-a\WeylQ_{ar}\,,
  \tag{25.4.2}\label{eq:25.4.2}
\end{equation}
so $\WeylQ_{(1/2)r}$ and $\WeylQ^\dagger_{(1/2)r}$ respectively lower and
raise the helicity by $1/2$.

We will first consider the case of simple supersymmetry. Consider a
supermultiplet with maximum helicity $\lambda_{\max}$, and let
$\ket{\lambda_{\max}}$ be any one-particle state with this helicity and
four-momentum $p^\mu$. Then
\begin{equation}
  \WeylQ^\dagger_{\frac12}\ket{\lambda_{\max}}=0\,,
  \tag{25.4.3}\label{eq:25.4.3}
\end{equation}
while acting on this state with $\WeylQ_{1/2}$ gives a state
$\ket{\lambda_{\max}-1/2}$ with helicity $\lambda_{\max}-1/2$. We will
define this state as
\begin{equation}
  \ket{\lambda_{\max}-1/2}
  \equiv(4E)^{-1/2}\WeylQ_{\frac12}\ket{\lambda_{\max}}\,.
  \tag{25.4.4}\label{eq:25.4.4}
\end{equation}
The fundamental anticommutation relation \eqref{eq:25.2.7} together with
Eqs.~\eqref{eq:25.4.1} and \eqref{eq:25.4.3} shows that this state is
normalized in the same way as $\ket{\lambda_{\max}}$
\begin{equation}
  \braket{\lambda_{\max}-1/2}{\lambda_{\max}-1/2}
  =\braket{\lambda_{\max}}{\lambda_{\max}}\,,
  \tag{25.4.5}\label{eq:25.4.5}
\end{equation}
and in particular this state cannot vanish. Eq.~\eqref{eq:25.2.32} shows that
$\WeylQ_{1/2}^2=0$, so acting with $\WeylQ_{1/2}$ on
$\ket{\lambda_{\max}-1/2}$ gives zero:
\begin{equation}
  \WeylQ_{\frac12}\ket{\lambda_{\max}-1/2}
  =(4E)^{-1/2}\WeylQ_{\frac12}^2\ket{\lambda_{\max}}
  =0\,.
  \tag{25.4.6}\label{eq:25.4.6}
\end{equation}
On the other hand, acting with $\WeylQ^\dagger_{1/2}$ on this state gives the
state with which we started. That is,
\[
  \begin{aligned}
  \WeylQ^\dagger_{\frac12}\ket{\lambda_{\max}-1/2}
  &=(4E)^{-1/2}
    \WeylQ^\dagger_{\frac12}\WeylQ_{\frac12}
    \ket{\lambda_{\max}}\\
  &=(4E)^{-1/2}
    \{\WeylQ^\dagger_{\frac12},\WeylQ_{\frac12}\}
    \ket{\lambda_{\max}}\,,
  \end{aligned}
\]
so that Eq.~\eqref{eq:25.4.1} and the anticommutation relation
\eqref{eq:25.2.31} yield
\begin{equation}
  \WeylQ^\dagger_{\frac12}\ket{\lambda_{\max}-1/2}
  =(4E)^{1/2}\ket{\lambda_{\max}}\,.
  \tag{25.4.7}\label{eq:25.4.7}
\end{equation}
Thus the supermultiplet consists of just two states, with helicities
$\lambda_{\max}$ and $\lambda_{\max}-1/2$. In the basis provided by these two
states, the operators $\WeylQ_{1/2}$ and $\WeylQ^\dagger_{1/2}$ are
represented by the matrices
\begin{equation}
  q_{\frac12}=\sqrt{4E}
  \begin{pmatrix}
    0&0\\
    1&0
  \end{pmatrix},
  \qquad
  q^\dagger_{\frac12}=\sqrt{4E}
  \begin{pmatrix}
    0&1\\
    0&0
  \end{pmatrix},
  \tag{25.4.8}\label{eq:25.4.8}
\end{equation}
while the operators $\WeylQ_{-1/2}$ and $\WeylQ^\dagger_{-1/2}$ are
represented by zero.

It is worth emphasizing that this is the \emph{only} kind of massless
supermultiplet in theories with simple supersymmetry. There are no massless
particles that are not accompanied with a superpartner, and none that have
more than one superpartner. Of course, CPT invariance implies that for every
supermultiplet of massless particles of helicities $\lambda$ and
$\lambda-1/2$ there must be an antimultiplet with helicities
$-\lambda+1/2$ and $-\lambda$. In particular, a massless particle and
antiparticle with helicities $+1/2$ and $-1/2$ must be accompanied by a
massless particle and antiparticle either with helicities $+1$ and $-1$ or
with helicities both zero.

How might the known quarks, leptons, and gauge bosons fit into this picture?
We will assume that the supersymmetry generator commutes with the generators
of the $SU(3)\times SU(2)\times U(1)$ gauge group\footnote{In simple
supersymmetry the generator $\WeylQ_\alpha$ must in any case commute with the
$SU(3)\times SU(2)$ generators, because semi-simple algebras like
$SU(3)\times SU(2)$ have no non-trivial one-dimensional representations.}.
The quarks and leptons belong to different representations of the gauge group
from the gauge bosons, so they cannot be in the same supermultiplets. We have
to conclude then that in the limit of high energy where $SU(2)\times U(1)$
symmetry breaking may be neglected, the massless quarks and leptons of each
color and flavor are in supermultiplets with \emph{pairs} of massless squarks
and sleptons of zero helicity and the same color and flavor, while the massless
gauge bosons are accompanied by massless gauginos of helicity $\pm1/2$
comprising an adjoint representation of $SU(3)\times SU(2)\times U(1)$.

Because gravity exists we know that in addition to the particles of the
standard model there must also exist a massless particle of helicity $\pm2$,
the \emph{graviton}. Massless particles with helicity $\lambda$ having
$|\lambda|>1/2$ must couple at low momentum to conserved
quantities.\footnote{This is discussed for the case of integer helicity in
Section 13.1. The argument for half-integer helicity was given by Grisaru and
Pendleton~\hyperref[ch25-ref-3]{[3]}.} Soft massless particles of helicity
$\pm1$ can couple to various internal symmetry generators, soft massless
particles of helicity $\pm3/2$ can couple to the supersymmetry generators
$\WeylQ_a$, and a soft particle of helicity $\pm2$ can couple to a single
conserved quantity, the momentum four-vector $P_\mu$, but there are no
conserved quantities to which a soft massless particle with $|\lambda|>2$
could couple. We conclude that the graviton cannot be in a supermultiplet with
particles of helicity $\pm5/2$, so it must be in a supermultiplet with a
massless particle of helicity $\pm3/2$, known as a \emph{gravitino}, coupled
to the supersymmetry generators themselves. The field theory of this
supermultiplet is known as \emph{supergravity}, and will be discussed in
Chapter 31.

Now let us consider the case of extended supersymmetry, with $N$
supersymmetry generators. We first note that because the
$\WeylQ_{(-1/2)r}$ all give zero when acting on the states of a
supermultiplet (including a state obtained by letting $\WeylQ_{(1/2)s}$ act on
any other state of the multiplet), the central charges $Z_{rs}$ must also
annihilate any state of the multiplet. With central charges out of the
picture, the supersymmetry generators $\WeylQ_{(1/2)r}$ all anticommute when
acting on a massless particle supermultiplet, so applying $n$ of them to a
one-particle state of maximum helicity $\lambda_{\max}$ and four-momentum
$p^\mu$ gives $N!/n!(N-n)!$ one-particle states of the same four-momentum
and helicity $\lambda_{\max}-n/2$, forming a rank $n$ antisymmetric tensor
representation of the $SU(N)$ $R$-symmetry\footnote{The $U(1)$ part of
$U(N)$ $R$-symmetry is often violated by quantum mechanical anomalies.}
\eqref{eq:25.2.30}. The maximum value of $n$ that gives a non-zero state is
$n=N$, so the minimum helicity in a supermultiplet is given by
\begin{equation}
  \lambda_{\min}=\lambda_{\max}-N/2\,.
  \tag{25.4.9}\label{eq:25.4.9}
\end{equation}
If we wish to exclude massless particle helicities $\lambda$ with
$|\lambda|>2$, then $\lambda_{\max}-\lambda_{\min}\leq4$, so extended
supersymmetries are allowed only with $N\leq8$.

For $N=8$, with helicities with $|\lambda|>2$ excluded, there is only one
possible supermultiplet, consisting of: 1 graviton with each helicity $\pm2$;
8 gravitinos with each helicity $\pm3/2$; 28 gauge bosons with each helicity
$\pm1$; 56 fermions with each helicity $\pm1/2$; and 70 bosons with helicity
zero.

Compare this with the case $N=7$, again excluding helicities with
$|\lambda|>2$. Here there are two supermultiplets. One supermultiplet
contains: 1 graviton with helicity $+2$; 7 gravitinos with helicity $+3/2$; 21
gauge bosons with helicity $+1$; 35 fermions with helicity $+1/2$; 35 bosons
with helicity zero; 21 fermions with helicity $-1/2$; 7 gauge bosons with
helicity $-1$; and 1 gravitino with helicity $-3/2$. The other is the
CPT-conjugate supermultiplet, with all helicities reversed. Adding the numbers
of particles in these two supermultiplets, we have 1 graviton with each
helicity $\pm2$; $7+1=8$ gravitinos with each helicity $\pm3/2$;
$21+7=28$ gauge bosons with each helicity $\pm1$; $35+21=56$ fermions with
each helicity $\pm1/2$; and $35+35=70$ bosons with helicity zero. Extended
supergravity theories with $N=8$ and $N=7$ thus have precisely the same
particle content and are in fact identical.

On the other hand, extended supergravity theories with $N\leq6$ have just $N$
gravitinos of each helicity $\pm3/2$ and are therefore all distinct.

For $N\leq4$ there is also the possibility of \emph{global} supersymmetry
theories, theories with supermultiplets that do not include gravitons or
gravitinos. For global $N=4$ supersymmetry there is just one supermultiplet,
containing: 1 gauge boson of each helicity $\pm1$; 4 fermions of each helicity
$\pm1/2$; and 6 bosons of helicity zero. This is equivalent to the global
supersymmetry theory with $N=3$, which has two supermultiplets: 1
supermultiplet with one gauge boson of helicity $+1$, 3 fermions with helicity
$+1/2$; 3 bosons with helicity zero; and 1 fermion with helicity $-1/2$; and
the other the CPT conjugate supermultiplet with opposite helicities. Adding
the numbers of particles of each helicity in these two $N=3$ supermultiplets
gives the same particle content as for $N=4$ global supersymmetry. The gauge
field theory with $N=4$ supersymmetry has remarkable properties, which will
be discussed in Section 27.9.

For $N=2$ extended global supersymmetry there are supermultiplets of two
different types, apart from those related by CPT. There are gauge
supermultiplets, each containing one gauge boson of helicity $+1$, two
fermions of helicity $+1/2$ forming a doublet under the $SU(2)$ $R$-symmetry
\eqref{eq:25.2.30}, and one boson of helicity zero, together with their
CPT-conjugate supermultiplets, with helicities reversed. Together each gauge
supermultiplet and its antimultiplet contain one gauge boson with each helicity
$\pm1$, an $SU(2)$ doublet of fermions of each helicity $\pm1/2$, and two
$SU(2)$ singlet bosons of helicity zero. Then there are
\emph{hypermultiplets}, which contain one fermion of each helicity $\pm1/2$
and an $SU(2)$ doublet of bosons of helicity zero, together with its
CPT-conjugate. (In quantum field theory a hypermultiplet cannot be its own
antimultiplet, because then the helicity zero particles would be described by
just two \emph{real} scalar fields, which cannot form an $SU(2)$ doublet.) Of
course, in the real world there would also have to be a graviton
supermultiplet, containing a graviton of helicity $+2$, an $SU(2)$ doublet of
gravitinos with helicity $+3/2$, and one gauge boson with helicity $+1$,
together with their CPT-conjugates, with opposite helicity. Gauge theories
with $N=2$ supersymmetry are constructed in Section 27.9, and explored
non-perturbatively in Section 29.5.

The particle content of these supermultiplets reveals a difficulty in
incorporating extended supersymmetry in realistic theories of particles at
accessible energies. In all cases but one, the helicity $+1/2$ fermions belong
to supermultiplets along with helicity $+1$ gauge bosons. Gauge bosons belong
to the adjoint representation of the gauge group, so if the supersymmetry
generators are invariant under the gauge group then the helicity $+1/2$
fermions must also belong to the adjoint representation, which is real. This
is in conflict with the fact that the known quarks and leptons belong to a
representation of $SU(3)\times SU(2)\times U(1)$ which is \emph{chiral}---that
is, for which the helicity $+1/2$ fermions belong to a complex representation,
which is then necessarily different from the representation furnished by
their CPT-conjugates, the helicity $-1/2$ fermions. The one exception, where
helicity $+1/2$ fermions are not in a supermultiplet with gauge bosons, is the
$N=2$ hypermultiplet discussed above. But in this case particles of both
helicities $+1/2$ and $-1/2$ are in the same supermultiplet, and therefore must
transform the same under any gauge transformations that leave the
supersymmetry generators invariant. They may belong to a complex
representation of this gauge group, but then the CPT-conjugate of this
hypermultiplet belongs to the complex-conjugate representation, and the
fermions of each helicity then belong to the sum of the two representations,
which is real, again in conflict with the chiral nature of the known quarks
and leptons.

In contrast, for simple supersymmetry there are supermultiplets containing
just helicity $+1/2$ and helicity zero, which may be in a complex
representation of the gauge group, distinct from the representation furnished
by the CPT-conjugate supermultiplet. Here there is no conflict with chirality.
For this reason, most discussions of supersymmetry as a symmetry that remains
unbroken at accessible energies have focused on simple rather than extended
supersymmetry.
